\section{Rozdział 1. Wstęp}
WSTĘP PROPOZYCJA (Motywacja → Przegląd źródeł → Metodologia → Metoda) [10-20 Stron] - (15\%-30\%) % usunąć
\subsection{Motywacja i uzasadnienie tematu}
% uzasadnieniem jest ciekawość, chęć pomocy przyszłym użytkownikom by podjąć decyzję które
% rozwiązanie wybrać, poprzez pokazanie konkretnych wyników na wykresach oraz umożliwienie przetestowania
% własnoręcznie wszystkich rozwiązań dla tych samych danych jednocześnie
\subsubsection*{Ważność problemu}
% tu pisze o tym, że ważne jest dla firm i użytkowników, by wiedzieć
% które rozwiązanie jest dokładniejsze, wydajniejsze, czy najtańsze
\subsubsection*{Znaczenie praktyczne i naukowe}
% tu możnaby napisać, że projekt jest gotowym rozwiązaniem do użytku w kontekście praktycznym oraz
% służy do porównania jakościowego oraz ilościowego modeli STT, co jest ważne dla nauki
\subsubsection*{Cel aplikacyjny/komercyjny}
% TODO:

\subsection{Problem badawczy i hipoteza}


\subsubsection*{Sformułowanie problemu}
\subsubsection*{Sformułowanie hipotezy}
\subsubsection*{Zakres badań}

\subsection{Przegląd literatury i stan wiedzy}
\subsubsection*{Przegląd literatury (min. 30 poz.)}
\subsubsection*{Aktualny stan wiedzy}
\subsubsection*{Analiza istniejących rozwiązań}
\subsubsection*{Identyfikacja luki badawczej}
\subsubsection*{Analiza rynku i konkurencji}

\subsection{Metodologia badań}
\subsubsection*{Tok postępowania badawczego}
\subsubsection*{Uzasadnienie wyboru metod}
\subsubsection*{Narzędzia i technologie}
\subsubsection*{Struktura pracy}
